\documentclass[12pt, letterpaper]{article}
\usepackage{graphicx} % Required for inserting images
\usepackage{polski}
%\usepackage[utf8]{inputenc}

\usepackage{hyperref}


\title{Programy użytkowe \\ Lista 1}
\author{ }
\date{}

\begin{document}

\maketitle

\begin{enumerate}
    \item Załóż konto w \href{https://www.overleaf.com/project}{Overleaf}. Utwórz pusty projekt.
    \begin{enumerate}
        \item Nadaj mu własny tytuł, imię i nazwisko autora oraz odpowiednią datę. Wprowadź tekst w języku polskim. Czy polskie znaki wyświetlają się poprawnie?
        \item Użyj polecenia \verb|\usepackage{polski}| w preambule  \\(tzn. nad 
        \verb|\begin{document}|), aby móc używać polskich znaków.
    \end{enumerate}
    \item Wpisz tekst podkreślony (\verb|\underline|{}), pochyły (\verb|\textit|{}) oraz pogrubiony \verb|\textbf{}|.
    \item Utwórz sekcję (\verb|\section{}|) oraz podsekcję(\verb|\subsection|).
    \item Wprowadź następujące formuły matematyczne (lub podobne):
    \begin{enumerate}
        \item $\cos(\pi)=-1$,
        \item $x^ax^b=x^{a+b}$,
        \item $\log_{10} 100=2,$
    \end{enumerate}
    używając odpowiednio: nazw funkcji i stałych (\verb|\cos,\pi|), znaków \\górnej(\verb|^|) oraz dolnej indeksacji(\verb|_|). W miejscach które tego wymagają, pamiętaj o nawiasach klamrowych (na przykład dla wykładnika $a+b$).
    
Użyj trzech trybów matematycznych:  \begin{itemize}
    \item \verb|$ ... $|,
    \item  \verb|$$ ... $$|,
    \item \verb|\begin{equation} ...\end{equation}|. 
\end{itemize}

Jakie dostrzegasz różnice w ich działaniu?

\item Utwórz \texttt{Example project} w Overleaf. 
\begin{enumerate}
    \item Zobacz, jak dodaje się grafikę. Podmień obrazek na własny.
    \item Dodaj polecenia, które pozwolą wprowadzić w nim polskie znaki.
\end{enumerate}
    
\end{enumerate}
\textbf{Zadanie domowe.} Wybierz dowolną formułę, prawo albo twierdzenie znane ze szkoły średniej (wymóg, to przynajmniej jedno równanie). Utwórz plik, w którym je zapiszesz. Tytuł pliku \textit{Zadanie domowe 1}. Dodaj swoje imię i nazwisko jako autora oraz datę.
\end{document}
